\section{Discussion}\label{sec:discussion}
\subsection{Interpretation of the Findings}
Under the combined attack, DIS-FDM and AEF-RPL flagged the same attackers, but downward delivery differed by up to 21 points and ARI by up to 0.31. Detection accuracy alone therefore did not tell how much service the network kept. When DIS-FDM isolates a flooder, the DIO storm stops, which removes the DIS part of the damage measured for H2. The false parents that the attacker advertised before it was isolated, however, remain in the root's routing table. AEF-RPL also rewrites those DAOs after detection, so both parts of the damage are removed. The ablation supports this reading: with repair disabled and detection left as it was, delivery fell to almost the unprotected level.

We introduced ARI for this kind of case. Conditional downward PDR considers only the packets that reached the root, whereas ARI also counts against the defense the nodes that lost their route and could no longer complete round trips. In our data the two metrics agreed in direction, but ARI showed a wider gap at high attacker fractions: with 14 static attackers it was 0.44 for DIS-FDM and 0.76 for AEF-RPL.

With learning, nodes flagged the same attackers as with the fixed rule, but earlier. A flood is contained sooner when its source is isolated sooner, and in the mobile networks the latency dropped by more than half. False isolations stayed at the same level. The learner works on the same evidence as the fixed rule and computes its reward from the same local indicators, so isolating an honest mobile node that exceeds the DIS threshold is rewarded. Reducing these isolations would require changing the evidence or the reward signal.

\subsection{Comparison with Prior Work}
Compared with DIS-FDM \cite{prajisha2025}, AEF-RPL keeps the DIS mitigation and its overhead reduction, adds detection of DAO-only attackers (100\% against 6.7\% with three static attackers), and restores downward service under the combined attack. It also supplies the accuracy and statistical evidence that DIS-FDM's evaluation lacks, including the window sensitivity that makes DIS-FDM miss slow floods. Li-MSD \cite{limsd2024} isolated far fewer honest nodes (0--2.7\% FPR) because it counts only DAOs and ignores DIS traffic, but it detected only 65.7\% of mobile attackers with seven attackers, missed almost all of them at slow rates, and preserved less service than AEF-RPL in every condition. Secure-RPL \cite{vermaranga2020} produced no false positives because it filters only early DIS arrivals, but it left the DAO damage untouched (77.4\% downward PDR with seven static attackers). The trade-off is thus between a defense that protects service on both planes (AEF-RPL) and defenses that isolate few honest nodes but protect only one plane (Li-MSD, Secure-RPL). Against QSec-RPL \cite{qsecrpl2023} and iTRPL \cite{itrpl2024}, which also learn, the comparison is architectural rather than numerical, because they target other attacks: AEF-RPL shows that a learner small enough for every constrained node (780 bytes) can speed up decisions, and that its accuracy is bounded by the evidence it learns from.

\subsection{Theoretical Implications}
Three points follow for theory. First, for insiders that attack more than one control plane, the response after detection affected the outcome as much as detection itself, and an evaluation limited to TPR and FPR would have rated DIS-FDM and AEF-RPL equally. Second, the Beta-calibrated uncertainty replaces hand-set belief masses with values estimated from each node's own calibration period, and Section~\ref{sec:theory} relates this uncertainty to the detection delay and the false-alarm bound. Third, when a learner has no ground truth, its reward carries over any bias present in the indicators used to compute it. At least one feedback signal that does not depend on the evidence is needed before learning can lower false positives.

\subsection{Practical Implications}
The results lead to two recommendations for operators of RPL meshes. The first-hop DAO-consistency check is inexpensive, since it uses a field the DAO already carries and needs neither keys nor additional messages. The ablation also shows that rewriting the false parent keeps honest descendants reachable, while dropping the DAO cuts them off. In networks with mobile nodes, the DIS threshold should take mobility into account. Moving honest nodes send DIS messages legitimately, and with the fixed counter rules more than 40\% of honest nodes were isolated in mobile runs, with DIS-FDM as well as with AEF-RPL. For protocol designers, the 780-byte learner is worth considering when reaction time matters, for instance with floods that drain batteries, as long as its false isolations are first measured in attack-free mobile runs.

\subsection{Contributions Revisited}
Each of the five contributions answers one or more of the gaps. AEF-RPL (contribution 1) covers both planes of a same-node attack, which is Gap~1, and estimates its evidence uncertainty from calibration data, which is Gap~2. The mathematical analysis (2) provides the delay, false-alarm, and value bounds that O6 asked for before any simulation. ARI and DRT (3) measure the service that remains, including route loss that delivery ratios do not capture. The paired evaluation with hypotheses fixed in advance (4) provides the accuracy and statistical evidence missing under Gap~3 and separates the effect of learning (Gap~4). Finally, the artifact (5) allows the numbers in this paper to be regenerated from the stored logs.

\section{Limitations and Threats to Validity}\label{sec:limitations}
The evaluation used 30-node networks. Attack-free pilots at 500 and 1000 nodes did not form stable bidirectional routes in the simulated stack, even after transport and timer adjustments, so the defense comparison was limited to the size at which the benign network was reliable. The pilots and their diagnostics are included in the artifact. The radio model is an ideal unit disk with CSMA and an always-on radio; a duty-cycled MAC or TSCH would change the cost of DIS flooding, as the H1 reproduction shows. Mobility followed the random waypoint model, and attackers followed fixed periodic schedules from a common onset. Detection is scored at the network level (a node counts as flagged if any observer flags it), and CPU cost was measured on the host, not on a microcontroller. Three items of the proposal's extended evaluation were not carried out: TSCH with Orchestra, intermittent DAO forgers, and memory measured on a mote platform. A short TSCH pilot without defenses is kept in the artifact, attack rate was varied instead of intermittency, and the memory cost is computed from the data structures (Section~\ref{sec:method}). The mobile false-isolation rate remains high for every method that uses DIS-FDM's DIS evidence, and the learner did not reduce it. The proposal also specified a per-node discount of the DIS and parent-change belief masses, estimated from attack-free warm-up windows. It was not implemented; instead, the DIS threshold is calibrated on each node and the parent-change thresholds are set per topology from attack-free runs. The blacklist logs indicate that the discount would not have changed the mobile false isolations, because 98\% of them in the attack-free mobile runs came from DIS-FDM's counter rules, which act on the counters directly and not on the belief masses. Ten seeds per cell limit statistical power, and the fresh-seed studies, which use broader multiplicity families, are therefore reported as replications of the pre-specified tests.

\section{Conclusion}\label{sec:conclusion}
This paper addressed RPL insiders that flood DIS messages and advertise false DAO parents at the same time, a combination that single-plane defenses detect but do not repair. We proposed AEF-RPL, which extends DIS-FDM with self-calibrated DAO-consistency evidence, persistence, selective DAO repair, and an on-node Q-learning decision layer, and we evaluated it in 2720 paired simulation runs.

AEF-RPL raised downward delivery and ARI over DIS-FDM in all six attacker-fraction and mobility conditions (98.3\% against 77.2\% delivery and ARI 0.92 against 0.71 with seven static attackers), detected the same attackers without a significant increase in false positives, and detected DAO-only attackers that DIS-FDM missed. Selective repair produced most of the service gain. On-node learning cut detection latency by 20--73\% but did not reduce false isolation, which originates in the DIS evidence of honest nodes that move. These results show that defenses against multi-plane routing attacks should be judged by the service they restore, not only by the attackers they flag.

Future work follows directly from the limitations: evaluating AEF-RPL in larger networks once stable bidirectional routing is established at that scale, on physical motes with duty-cycled MACs, against intermittent and adaptive attackers, and with a DIS threshold that adapts to mobility and a learning reward that includes a feedback signal independent of the DIS evidence, so that learning can reduce false isolation under mobility.

\section*{Data and Code Availability}
The firmware, scenario generators, all 2720 raw simulation logs with checksums, and the analysis scripts are available at \url{https://github.com/KarimCh01/rpl-dis-dao-evidence-fusion}. The command \texttt{make reproduce-paper} regenerates every table, statistical test, and figure in this paper from the stored logs without running the simulator.

\section*{Ethics Statement}
The study used only simulated networks and synthetic traffic. No human participants, personal data, or real networks were involved.


\section*{Funding and Competing Interests}
This research received no specific grant from any funding agency. The author declares no competing interests.

\section*{Acknowledgment}
An AI-based writing assistant was used only to check the grammar and wording of this manuscript. It did not generate the technical content, experiments, data, analysis, or conclusions. The author reviewed every suggested change and takes full responsibility for the text.
